% chapters/preface.tex — 머리말

\chapter*{머리말}
\addcontentsline{toc}{chapter}{머리말}
\markboth{머리말}{머리말}

\section*{왜 이 책을 썼는가}

2022년 말 ChatGPT가 공개된 이후, 대형 언어 모델(LLM)은 인공지능 분야의 가장 뜨거운 주제가 되었다. 수많은 개발자와 연구자가 HuggingFace Transformers를 통해 이미 만들어진 모델을 다운로드하고, 파인튜닝하고, 서비스에 통합하는 방법을 빠르게 익혔다. 하지만 ``이 모델이 내부적으로 어떻게 동작하는가?''라는 질문에 자신 있게 답할 수 있는 사람은 여전히 소수다.

이 책은 그 질문에 답하기 위해 쓰였다. 토크나이저를 직접 만들고, 어텐션을 손으로 구현하고, 트랜스포머 블록을 조립하고, 학습 루프를 돌려보면서 LLM의 뼈대를 이해하는 것이 이 책의 목표다. 추상화된 라이브러리를 쓰지 않고 PyTorch의 기본 연산만으로 처음부터 끝까지 직접 만든다. 그 과정에서 겪는 모든 디테일이 결국 LLM을 ``이해''하는 길이다.

\section*{이 책의 특징}

이 책은 단순히 이론을 나열하지 않는다. \textbf{코드를 먼저 만들고, 단위 테스트로 검증한 뒤, 그 과정을 책으로 풀어낸다.} 책에 등장하는 모든 Python 코드는 실제로 실행되며, pytest 기반 단위 테스트를 통과한 코드만 수록되었다. 독자는 GitHub 저장소에서 전체 코드를 다운로드하여 직접 실행하고 수정하며 학습할 수 있다.

각 장은 이전 장에서 구현한 모듈을 import하여 사용한다. 예를 들어 4장의 임베딩 모듈은 3장의 토크나이저를 사용하고, 5장의 어텐션은 4장의 임베딩을 사용한다. 이 ``쌓아가는 구조''를 통해 독자는 LLM이 어떻게 조립되는지를 단계적으로 체감하게 된다. 수학 수식은 직관적 설명 다음에 등장하며, 초보자도 흐름을 놓치지 않도록 배려했다.

\section*{누가 읽어야 하는가}

이 책의 주 독자는 Python 기초 문법을 알고 있으나 딥러닝은 처음인 개발자다. HuggingFace Transformers를 써본 적은 있으나 내부 동작을 모르는 엔지니어, ``Attention Is All You Need'' 논문을 읽어보았으나 코드로 구현해본 적은 없는 연구자도 포함된다. 미적분과 기초 선형대수를 가정하되, 행렬 미분 등의 고급 주제는 부록에서 보충 설명한다.

\begin{flushright}
\textit{2026년 여름}\\[0.3em]
\textbf{dirmich}
\end{flushright}

\newpage

\section*{감사의 글}

이 책을 쓰는 동안 수많은 오픈소스 프로젝트와 논문의 도움을 받았다. 특히 Andrej Karpathy의 nanoGPT, HuggingFace Tokenizers, 그리고 PyTorch 개발팀의 훌륭한 문서에 깊은 감사를 전한다. 이들은 LLM의 민주화에 헌신한 선구자들이다.

또한 이 책의 코드가 기반을 둔 수많은 연구자들에게 감사드린다. Ashish Vaswani와 트랜스포머 원논문의 공동 저자들은 이 모든 것의 출발점을 만들어주었다. Alec Radford와 OpenAI 팀은 GPT 시리즈를 통해 스케일링의 가능성을 보여주었다. Tianyi Zhang과 Jason Wei는 instruction tuning과 정렬의 기반을 닦았다.

가족과 친구들에게도 감사를 전한다. 글을 쓰는 동안 묵묵히 기다려준 이들의 인내가 없었다면 이 책은 완성되지 않았을 것이다.

마지막으로, 이 책을 읽는 모든 독자에게. 여러분의 호기심과 열정이 LLM의 미래를 만든다. 직접 만들고, 실험하고, 실패하고, 다시 시도하라. 그것이 학습의 본질이다.

\begin{flushright}
\textbf{dirmich}\\
\textit{Highmaru Press}
\end{flushright}
